\begin{onehalfspace}

Questa sezione documenta per esteso l'analisi M8/M13 sintetizzata nella
Sezione~\ref{sec:risultati_logico} del Capitolo~\ref{chap:risultati}.
Conserva la separazione fra livelli, la griglia completa, l'analisi delle
ripetizioni e il contesto fault-tolerant; nel corpo rimangono definizione
del proxy, curva principale e conclusione.

\subsection{Separazione dei quattro livelli}
\label{app:shor_quattro_livelli}

\begin{table}[H]
\centering
\footnotesize
\begin{tabular}{clll}
\toprule
\textbf{Livello} & \textbf{Oggetto} & \textbf{Strumento} & \textbf{Domanda} \\
\midrule
1 & Shor ideale & Qiskit Aer & limite senza rumore \\
2 & Shor rumoroso & Aer + noise model & sensibilità agli errori fisici \\
3 & QEC isolata & Qiskit; Stim/PyMatching & soppressione dell'errore \\
4 & proxy per gate & Aer + canale Pauli & sensibilità a un residuo uniforme \\
\bottomrule
\end{tabular}
\caption[Architettura a quattro livelli della parte sperimentale]{I
quattro livelli isolano domande diverse. Il quarto non è una simulazione
fault-tolerant e non riceve automaticamente i tassi aggregati del terzo.}
\label{tab:quattro_livelli}
\end{table}

Una simulazione esplicita incontra una barriera dimensionale --- sostituire
ognuno dei dodici qubit di $N=15$ con un blocco codificato conduce a
centinaia di qubit --- e una strutturale: Stim è efficiente sui circuiti
stabilizer, mentre aritmetica modulare e QFT rendono Shor non-Clifford.
Aer simula il circuito generale piccolo e Stim/PyMatching caratterizza la
QEC isolata; nessuno dei due equivale a una compilazione fault-tolerant.

Il manifest M8 registra 224 CX, 302 RZ, 70 SX, 8 misure e lo SHA-256 con
prefisso \texttt{c7f4cf3b} e suffisso \texttt{a50f4ba2c}. Ogni punto aggrega 20 repliche da
4096 shot, cioè $81\,920$ misure. Il successo per misura vale uno quando
frazioni continue e controlli MCD restituiscono fattori non banali.

\subsection{Contesto delle stime fault-tolerant}
\label{app:shor_risorse}

Le stime crittografiche comprendono routing, cicli di sindrome, fabbriche
di stati magici e schedulazione; non derivano dall'estrapolazione di
$N=15$. Gidney ed Eker\aa{} stimavano circa venti milioni di qubit e otto
ore per RSA-2048~\cite{gidney2021}; lavori successivi riducono il costo
grazie ad aritmetica, codifica e preparazione delle risorse, ma sotto
assunzioni architetturali differenti.

\begin{table}[H]
\centering
\small
\begin{tabular}{p{0.26\textwidth}p{0.08\textwidth}p{0.20\textwidth}p{0.30\textwidth}}
\toprule
\textbf{Lavoro} & \textbf{Anno} & \textbf{Bersaglio} & \textbf{Qubit fisici e tempo} \\
\midrule
Gidney ed Eker\aa{}~\cite{gidney2021} & 2019 & RSA-2048 & $\sim$20 milioni, 8 ore \\
Gidney~\cite{gidney2025} & 2025 & RSA-2048 & $<$1 milione, $<$1 settimana \\
Pinnacle, qLDPC~\cite{pinnacle2026} & 2026 & RSA-2048 & $<$100\,000 (preprint) \\
Cain \emph{et al.}~\cite{cain2026} & 2026 & P-256; RSA-2048 & $\sim$10\,000 atomici; P-256 in giorni con 26\,000 (preprint) \\
\bottomrule
\end{tabular}
\caption[Traiettoria delle stime di risorsa per Shor su scala crittografica]{Le
righe non sono direttamente confrontabili: cambiano codice, modello
d'errore, ciclo e criterio di successo. Le due voci 2026 sono preprint e
nessun valore discende dalle curve di questa tesi.}
\label{tab:traiettoria_risorse}
\end{table}

La discesa delle stime è trainata soprattutto da correzione e
compilazione, ma resta una valutazione prospettica. La curva M8 è una
sensibilità locale e non una previsione di risorse.

\subsection{Griglia M8 completa e probabilità cumulativa}
\label{app:shor_m8_dettaglio}

La griglia completa con intervalli Wilson, l'Equazione~\ref{eq:cumulativa} e
la curva delle ripetizioni sono mantenute nel corpo, rispettivamente nella
Tabella~\ref{tab:m8_griglia_completa} e nella
Figura~\ref{fig:qec_shor_pk}. Qui resta il contesto che impedisce di tradurre
automaticamente quei punti in requisiti fault-tolerant.

\subsection{Perché le metriche QEC non entrano direttamente in M8}
\label{app:shor_metriche}

\begin{table}[H]
\centering
\small
\begin{tabular}{p{0.22\textwidth}p{0.38\textwidth}p{0.26\textwidth}}
\toprule
\textbf{Grandezza} & \textbf{Unità/contesto} & \textbf{Inseribile direttamente in M8?} \\
\midrule
$p_L$ di M8 & Pauli non-identità dopo ogni SX/X/CX & sì, per definizione \\
errore Steane & fallimento del protocollo protetto & no \\
errore surface code & fallimento logico per memoria/ciclo e distanza & no \\
guadagno decoder & differenza nel benchmark correlato & no \\
\bottomrule
\end{tabular}
\caption[Contratto tra risultati QEC e curva M8]{La conversione richiede
circuito logico, cicli, operazioni fault-tolerant e correlazioni residue.}
\label{tab:decoder_su_shor}
\end{table}

Una verifica recente su 114 atomi mostra che una variante pre-compilata
di Shor su qubit logici dimezza fino a un fattore due la distanza in
variazione totale rispetto all'esecuzione fisica; ladder CX logiche
mostrano errori due--quattro volte inferiori e un codice $[[16,4,4]]$ un
fattore otto~\cite{rines2025}. Il risultato riguarda il vantaggio del
livello logico, non una fattorizzazione scalabile e resta soggetto alla
critica di Smolin~\cite{smolin2013}. La cifra di 48 qubit logici su 280
atomi appartiene invece al distinto esperimento di Bluvstein
\emph{et al.}~\cite{bluvstein2024}, non a Shor. Questa evidenza conferma
solo la direzione qualitativa del proxy; non fornisce la conversione
mancante.

\end{onehalfspace}
